\documentclass[10pt,a4paper]{amsart}
\usepackage[margin=19mm]{geometry}
\usepackage{amsmath,amssymb,mathtools}
\usepackage[scr=rsfs,cal=euler]{mathalfa}
\usepackage{xfrac}
\usepackage[unicode,hidelinks,bookmarks=false]{hyperref}
\newcommand{\ie}{i.e.\ }
\newcommand{\cf}{cf.\ }
\newcommand{\resp}{respectively\ }
\newcommand{\Subsection}{\subsection}
\providecommand{\nobreakdash}{\nobreak\hbox{-}}
\setlength{\emergencystretch}{2em}
\multlinegap0pt
\usepackage{bm}
\usepackage{bbm}
\usepackage{dsfont}
\usepackage{xspace}
\usepackage{cancel}
\usepackage{mathtools}
\usepackage{amsthm}
\makeatletter
\@ifundefined{lem}{\newtheorem{lem}{Lem}}{}
\@ifundefined{prop}{\newtheorem{prop}[lem]{Proposition}}{}
\makeatother
\newcommand{\leg}{\overwithdelims ()}
\newcommand{\new}{\mathrm{new}}
\renewcommand{\mod}{\bmod}
\title[Distribution of local signs of modular forms and murmuration]{Distribution of local signs of modular forms and murmurations of Fourier coefficients}
\author{Kimball Martin}
\date{Source: arXiv:2409.02338v3 (CC BY 4.0)}
\begin{document}
\maketitle

\noindent\emph{Source and licence.} Distribution of local signs of modular forms and murmurations of Fourier coefficients. Version arXiv:2409.02338v3 (CC BY 4.0), distributed under the Creative Commons Attribution 4.0 International licence, as recorded in the arXiv licence metadata for this exact version and shown on the abstract page https://arxiv.org/abs/2409.02338v3. Full rights notice: licenses/arxiv-2409.02338-CC-BY-4.0.txt.\par\smallskip

\noindent\textit{Editorial scope.} Counted unit: the Prop whose source label is \texttt{prop:req1} in the article (source0.tex lines 934--956, source label \texttt{prop:req1}), delivered together with the article's own complete original proof of it (source0.tex lines 958--986, one \texttt{proof} environment of 29 lines). No mathematics is written, repaired or re-derived: statement and proof are the article's own text line for line. Required context retained uncounted: none. Every definition, display and label of those lines is retained unchanged. Omitted from this package and not redistributed: 1--933, 957--957, 987--1677. Nothing inside a retained excerpt is omitted or altered: every retained body is delivered exactly as the article's lines are written, so no omission receipt of any kind accompanies this package. Citations: the retained excerpts carry no citation command at all, so no bibliography entry of the article is transcribed and no reference is redistributed. Reference adaptation: none was needed and none was made; every reference inside the delivered unit resolves inside it, so no unresolved reference or citation is produced. Printed slips kept literal: no printed word, symbol, hypothesis or number of the counted statement or of its proof was altered. Layout and class: the article's own class is \texttt{amsart} with packages including amsmath, amsthm, amscd, amssymb, color, thmtools, hyperref, amsrefs, cleveref, enumerate, geometry, graphicx, caption, array; the wrapper is the lane's pinned \texttt{amsart} editorial wrapper at 10pt on A4, which restates the theorem-like environments the retained text needs and adds the article's own macro definitions that the retained text needs (\texttt{\textbackslash leg}, \texttt{\textbackslash mod}, \texttt{\textbackslash new}), with extra wrapper packages (bm, bbm, dsfont, xspace, cancel, mathtools). The delivered numbering is the unit's own. Assets: no retained span contains a figure environment, a graphics-inclusion command, a PDF-page inclusion, a table or a TikZ picture; no image file is copied into this package, the package declares no asset dependency and no image file is redistributed with it. Licence: version arXiv:2409.02338v3 of this article is distributed under the Creative Commons Attribution 4.0 International licence, as recorded in the arXiv licence metadata for this exact version and shown on the abstract page https://arxiv.org/abs/2409.02338v3. A full rights notice is delivered at licenses/arxiv-2409.02338-CC-BY-4.0.txt. Characters: every retained excerpt is pure ASCII, so the delivered engine, XeLaTeX with its default Latin Modern text font, reports no missing character.
\begin{prop} \label{prop:req1}
Let $q$ be a prime, $M \ge 1$ be coprime to $q$.  Write $M=2^e M'$, where $M'$ is odd.

\begin{enumerate}
\item Suppose $q \ge 5$ and either $k \ge 4$ or $M$ is not squarefree.  Then
$\Delta_k(q,M) = 0$ if and only if 
(i) $M'$ is not cubefree; (ii) ${-q \leg p} = 1$ for some $p \parallel M'$; 
(iii) $16 \mid M$ and $q \equiv 1 \mod 4$; (iv) $32 \mid M$ and $q \equiv 3 \mod 8$;
or (v) $v_2(M) \ne 0, 4$ and $q \equiv 7 \mod 8$.

\item Suppose $q \ge 5$, $k=2$ and $M$ is squarefree.
Then $\Delta_k(q,M) = 0$ if and only if (i) $M = 1$ and
$q \in \{ 5, 7, 13, 17 \}$; or (ii) $M = 2$ and $q \in \{ 5, 11, 13, 19, 37, 43, 67, 163 \}$.

\item Suppose $q=2$.  We have $\Delta_k(2,M) = 0$ if and only if 
(i) $M$ is not cubefree; (ii) $\prod_{p^2 \parallel M} {2 \leg p}$ is $-1$
if $k \equiv 0, 2 \mod 8$ and $+1$ if $k \equiv 4, 6 \mod 8$; or
(iii) $k=2$ and $M = 1$ or $M \equiv 3, 5 \mod 8$ is prime.

\item Suppose $q=3$.  If $k \ge 4$ or $M$ is not squarefree, then  $\Delta_k(3,M) = 0$ if and only if (i) $M'$ is not cubefree; (ii) ${-3 \leg p} = 1$ for some $p \parallel M'$; (iii) $32 \mid M$; (iv) $M$ is odd and $k \equiv 4, 10 \mod 12$; or (v) $4 \parallel M$ and $k \not \equiv 4, 10 \mod 12$.
 If $k=2$ and $M$ is squarefree, then $\dim S_k^\new(3M)^{+_3} = \dim S_k^\new(3M)^{-_3}$ if and only if $M = 1, 2$. 
\end{enumerate}
\end{prop}

\begin{proof} 

Case (1) follows immediately from the above vanishing conditions for $\alpha_1(-q;e)$.

Case (2): Now suppose $q \ge 5$, $k=2$ and $M$ is squarefree.  Then $\Delta_k(q,M) = 0$ if and only if $\frac 12 \alpha_1(-q;e) \kappa_{-q}(M') = \mu(M)$.  

For a discriminant $\Delta < -4$, $H(\Delta) \ge 1$.
Also if $\Delta = \lambda^2 \Delta_0$ where $\Delta_0$ is a fundamental discriminant and $\lambda > 1$, then $H(\Delta) \ge 2 H(\Delta_0)$.  Thus the only integers $\Delta$ such that $H(\Delta) = 1$ are $\Delta = \{ -7, -8, -11, -19, -43, -67, -163 \}$.  Similarly, there are 19 discriminants $\Delta < 0$ with $H(\Delta) = 2$ (the minimal one is $\Delta = -427$).

If $\kappa_{-q}(M') \ne 0$, it equals $(-2)^{\omega(M')}$.  Thus we can only have
$\Delta_k(q,M) = 0$ if (i) $M' = 1$ and $\alpha_1(-q; e) = \pm 2$, or (ii)  $M' = p$ is prime and $\alpha_1(-q; e) = \pm 1$.  

When $M=1$, $\Delta_k(q,M) = 0$ if and only if $H(-4q) = 2$, i.e., if $q \in \{ 5, 7, 13, 17 \}$.  When $M = 2$, $\Delta_k(q,M) = 0$ if and only if $H(-4q) = 2 H(-q) + 2$, i.e., if and only if $q \in \{ 5, 11, 13, 19, 37, 43, 67, 163 \}$.
When $M = p \ge 3$, $\Delta_k(q,M) = 0$ if and only if $H(-4q) = 1$ and ${-q \leg p} = -1$, which never happens.
When $M = 2p$, for a prime $p \ge 3$, $\Delta_k(q,M) = 0$ if and only if $2H(-q) = H(-4q) + 1$, which also never happens.  This finishes case (2).

Case (3): Next suppose $q=2$, and $k \ne 2$ or $M$ is not squarefree.  Then
$t^\new(1;M) = 0$ if and only if $\kappa_{-2}(M) = \kappa_{-1}(M) = 0$
or $(-1)^{k/2} p_k(\sqrt 2,1) = \prod_{p^2 \parallel M} {2 \leg p}$.  The former never happens.  The latter condition means $\prod_{p^2 \parallel M} {2 \leg p}$ is $-1$ if $k \equiv 0, 2 \mod 8$ and $+1$ if $k \equiv 4, 6 \mod 8$.

If $q=2$, $k=2$ and $M$ is squarefree, then one needs $\kappa_{-2}(M) + \kappa_{-1}(M) = 2\mu(M)$, which is true when $M = 1$ or $M$ is a prime $p \equiv 3, 5 \mod 8$.  This proves case (3).

Case (4): Finally suppose $q=3$.  Then $\Delta_k(q,M) = c \kappa_{-3}(M') + \delta_{k=2} \mu(M)$, where $c = \frac 12(-1)^{k/2} \alpha_1(-3;e) - \frac 13 p_k(\sqrt 3,1) \kappa_{-3}(2^e)$.  One checks that $c \in \{ -1, 0, 1 \}$, and
$c = 0$ if and only if (i) $e \ge 5$; (ii) $e =  0$ and $k \equiv 4, 10 \mod 12$; or (iii) $e = 2$ and $k \equiv 0, 2, 6, 8 \mod 12$.

Hence if $k \ne 2$ or $M$ is not squarefree, then $\Delta_k(q,M) = 0$ if and only if one of (i)--(iii) above holds; (iv) $M'$ is not cubefree; or (v) ${-3 \leg p} = 1$ for some odd $p \parallel M$.

Now suppose $k=2$ and $M$ is squarefree.  Then $c = (-1)^{e+1}$, and we see $\mu(M) = c \kappa_{-3}(M')$ if and only if $M = 1, 2$.  This completes case (4).
\end{proof}

\end{document}
